\chaptertitle{参考答案}

\sectiontitle{第一章}

\textbf{第1讲}

1.(a) $2n-1$ \quad (b) $2n$ \quad (c) $3n-1$ \quad (d) $3n-2$\\
2.(a) $a+b$ \quad (b) $3x$ \quad (c) $mn$ \quad (d) $a^2+b^3$ \quad (e) $\dfrac{x}{2}+2y$ \quad (f) $5(a-b)$\\
3.(a) $3a$ \quad (b) $abc$ \quad (c) $5x+2y$ \quad (d) $3a+4b$\\
4. $2n+1$根\\
5.(a) $ax$ \quad (b) $xa$ \quad (c) 相等，乘法交换律

\textbf{第2讲}

1.(a) $2x-5$ \quad (b) $3(a+b)$ \quad (c) $m^2-n$ \quad (d) $n-1$和$n+1$ \quad (e) $\dfrac13x=\dfrac12x-4$ \quad (f) $2(a+(a-3))=4a-6$\\
2.(a) 11 \quad (b) 5 \quad (c) 16 \quad (d) 9\\
3.(a) 4 \quad (b) 3 \quad (c) 1 \quad (d) 15\\
4.(a) $8x$ \quad (b) $5a$ \quad (c) $2m+4n$ \quad (d) $x^2+2x+5$ \quad (e) $8ab-3a^2$ \quad (f) $x^2+3x-2$\\
5.(a) 错，$3x+2y$不能合并 \quad (b) 错，$5a^2$ \quad (c) 错，$3x$ \quad (d) 对

\textbf{第3讲}

1.(a) $4x-y$ \quad (b) $a^2-5ab+7b^2$ \quad (c) $n$ \quad (d) $5x^2+2xy-3$ \quad (e) $-2x+2y$ \quad (f) $2a+3b$\\
2.(a) $2x^2+2$ \quad (b) $2x^2-5x+7$ \quad (c) $-4a^2-5ab$ \quad (d) $2x^2-y^2$\\
3.(a) $4x+y$ \quad (b) $2x+3y$ \quad (c) $3x+7y$\\
4. $(3x^2+x-2)-(2x^2-3x+1)=x^2+4x-3$\\
5. $2((2a+3)+(a-1))=2(3a+2)=6a+4$

\sectiontitle{第二章}

\textbf{第4讲}

1.(a) $x=8$ \quad (b) $x=40$ \quad (c) $x=9$ \quad (d) $x=72$ \quad (e) $x=4$ \quad (f) $x=6$\\
2.(a) $x=9$ \quad (b) $x=8$ \quad (c) $x=3$ \quad (d) $x=25$ \quad (e) $x=7$ \quad (f) $x=16$\\
3.(a) $3x+5=26$，$x=7$ \quad (b) $\dfrac{x}{2}-7=3$，$x=20$ \quad (c) $4x-25=7$，$x=8$\\
4.(a) 对 \quad (b) 对 \quad (c) 错，$x=16$

\textbf{第5讲}

1.(a) $x=6$ \quad (b) $x=5$ \quad (c) $x=6$ \quad (d) $x=6$\\
2.(a) $x=2$ \quad (b) $\dfrac{17}{5}$ \quad (c) $\dfrac92$ \quad (d) $-\dfrac15$\\
3.(a) $x=30$ \quad (b) $x=3$ \quad (c) $x=22$ \quad (d) $x=1$\\
4.(a) 30 \quad (b) 甲100本，乙40本 \quad (c) 甲2100元，乙1500元 \quad (d) 甲17岁，乙11岁

\textbf{第6讲}

1.(a) $x<3$ \quad (b) $x>5$ \quad (c) $x\le2$ \quad (d) $x\ge4$\\
2.(a) $x<1$ \quad (b) $x\ge-1$ \quad (c) $x\le1$ \quad (d) $x\le3$\\
3.(a) $3x-5\ge10$ \quad (b) $x-25\ge15$ \quad (c) $\dfrac{x-40}{40}\ge20\%$\\
4.(a) $x=3$ \quad (b) $x>3$ \quad (c) $x\le0$\\
5. 提示：$-3<-1<1<3$。乘$-1$后数轴上的位置反过来。

\sectiontitle{第三章}

\textbf{第7讲}

1.(a) $\begin{cases}x=5\\y=2\end{cases}$ \quad (b) $\begin{cases}x=\dfrac{22}{7}\\y=\dfrac{17}{7}\end{cases}$ \quad (c) $\begin{cases}x=3\\y=1\end{cases}$\\
2.(a) $\begin{cases}x=3\\y=2\end{cases}$ \quad (b) $\begin{cases}x=5\\y=0\end{cases}$ \quad (c) $\begin{cases}x=2\\y=-1\end{cases}$\\
3.(a) 甲20，乙12 \quad (b) 篮球8个，排球4个 \quad (c) 客车6辆，货车6辆\\
4. 无穷多组解（两个方程实质相同，代表同一条直线）

\textbf{第8讲}

1.(a) $-3<x<3$ \quad (b) $2<x<4$ \quad (c) $x\ge2$ \quad (d) $-1<x<2$\\
2. $a>0$时$x<\dfrac ba$；$a<0$时$x>\dfrac ba$；$a=0$且$b>0$时全体实数，$a=0$且$b\le0$时无解\\
3. $50\le$售价$\le52$元

\textbf{第9讲}

1.(a) $x=2$ \quad (b) $x>2$ \quad (c) $x<1$ \quad (d) $y=2$\\
2.(a) $x=3$ \quad (b) $x\le3$ \quad (c) $y=4$ \quad (d) $x=\dfrac32$\\
3.(a) $x=-2,-1,0,1,2$对应$y=-8,-5,-2,1,4$ \quad (b) 3 \quad (c) $k$（斜率）\\
4.(a) 对 \quad (b) 对（斜率相同，截距不同） \quad (c) 错（$b=1$，不过原点）\\
5. $y=2x+1$是函数关系（一条直线）；$2x+1=0$是方程（一个点的$x$坐标）

\sectiontitle{第四章}

\textbf{第10讲}

1.(a) $2x^2-7x-4$ \quad (b) $6a^2+5ab-6b^2$ \quad (c) $x^2-xy-6y^2$ \quad (d) $6x^3+x^2-4x+1$\\
2.(a) $x^2-36$ \quad (b) $9a^2-4b^2$ \quad (c) $25m^2-n^2$ \quad (d) $a^4-b^2$ \quad (e) 9991 \quad (f) 39996\\
3.(a) $x^2+10x+25$ \quad (b) $4a^2-12a+9$ \quad (c) $9m^2+12mn+4n^2$ \quad (d) 10404 \quad (e) 9801 \quad (f) 40401\\
4.(a) 错，$a^2+2ab+b^2$ \quad (b) 错，$a^2-2ab+b^2$ \quad (c) 对 \quad (d) 对\\
5. $(x+y)^2-(x-y)^2=4xy$

\textbf{第11讲}

1.(a) $5(x+2)$ \quad (b) $3a(a-2)$ \quad (c) $3x^2(2x-3)$ \quad (d) $4mn(m-2n)$ \quad (e) $(x-3)(2x+5)$ \quad (f) $(x+y)(3a-2b)$\\
2.(a) $(x+4)(x-4)$ \quad (b) $(2a+3)(2a-3)$ \quad (c) $(5m+n)(5m-n)$ \quad (d) $(3x+4y)(3x-4y)$\\
3.(a) $(x+3)^2$ \quad (b) $(2a-5)^2$ \quad (c) $(3m+2n)^2$ \quad (d) $(x-4)^2$\\
4.(a) $(x+3)(x+4)$ \quad (b) $(x-5)(x+1)$ \quad (c) $(x+5)(x-2)$ \quad (d) $(x-3)(x-4)$\\
5.(a) $3(x+3)(x-3)$ \quad (b) $2(x-2)^2$ \quad (c) $x(x+2)(x-2)$ \quad (d) $y(x+1)(x-1)$

\textbf{第12讲}

1.(a) $2x$ \quad (b) $\dfrac{2a}{b}$ \quad (c) $x-1$ \quad (d) $x+2$\\
2.(a) $\dfrac{a+b}{ab}$ \quad (b) $\dfrac{2x-3}{x^2}$ \quad (c) $\dfrac{2x}{x^2-1}$ \quad (d) $\dfrac{a^2-b^2}{ab}$\\
3.(a) 4 \quad (b) $\dfrac{y}{x}$ \quad (c) $a+1$\\
4.(a) $3a\sqrt2$ \quad (b) $2x\sqrt{2xy}$ \quad (c) $5ab^2\sqrt2$\\
5.(a) $\dfrac{\sqrt3}{3}$ \quad (b) $\dfrac{2\sqrt5}{5}$ \quad (c) $\dfrac{\sqrt3-1}{2}$

\sectiontitle{第五章}

\textbf{第13讲}

1.(a) $x_1=0,x_2=3$ \quad (b) $x_1=-3,x_2=-4$ \quad (c) $x_1=-1,x_2=5$ \quad (d) $x_1=3,x_2=-3$ \quad (e) $x_1=x_2=-3$ \quad (f) $x_1=7,x_2=-7$\\
2.(a) $x_1=3,x_2=-1$ \quad (b) $x_1=4,x_2=-2$ \quad (c) $x_1=5,x_2=1$ \quad (d) $x_1=-1,x_2=\dfrac52$\\
3.(a) 0或1 \quad (b) 7和8或-7和-8 \quad (c) 长9米，宽5米

\textbf{第14讲}

1.(a) $x_1=1,x_2=-5$ \quad (b) $x_1=4,x_2=2$ \quad (c) $x_1=9,x_2=-1$ \quad (d) $x_1=-2,x_2=-3$\\
2.(a) $\dfrac{5\pm\sqrt{13}}{2}$ \quad (b) $x_1=3,x_2=\dfrac12$ \quad (c) $x_1=\dfrac13,x_2=-2$ \quad (d) $2\pm\sqrt2$\\
3.(a) $\Delta=1>0$，两个不等实根 \quad (b) $\Delta=-3<0$，无实根 \quad (c) $\Delta=0$，两个相等实根 \quad (d) $\Delta=17>0$，两个不等实根\\
4.(a) 5 \quad (b) 6 \quad (c) $\dfrac56$ \quad (d) 13\\
5. 能因式分解就用因式分解（最快）；不能分解但系数简单就用配方；通用用公式

\textbf{第15讲}

1.(a) $x_1=3,x_2=-1$ \quad (b) $x<-1$或$x>3$ \quad (c) $-1<x<3$ \quad (d) 顶点$(1,-4)$\\
2.(a) $x_1=3,x_2=1$ \quad (b) $1<x<3$ \quad (c) 开口向下，顶点$(2,1)$\\
3.(a) $x$轴，横坐标 \quad (b) 2 \quad (c) 相切（1个交点） \quad (d) 没有交点\\
4.(a) $y=(x-3)^2-1$ \quad (b) 顶点$(3,-1)$，对称轴$x=3$ \quad (c) $x<2$或$x>4$\\
5. 相同点：都是问"等于多少""大于/小于多少""每个$x$对应什么$y$"。不同点：一次是直线（最多1个解），二次是抛物线（最多2个解）；一次不等式是"一侧"，二次可能是"中间或两边"。

\sectiontitle{第六章}

\textbf{第16讲}

1.(a) $x=3$ \quad (b) 无解 \quad (c) $x=3$ \quad (d) $x=\dfrac74$\\
2.(a) $x=5$ \quad (b) $x=3$ \quad (c) $x=3$ \quad (d) $x=4$\\
3. 不一样。分式方程因为去分母时可能乘以0；根式方程因为平方时正负信息丢失。

\textbf{第17讲}

1.(a) 恒成立（任意实数） \quad (b) $x=\dfrac{21}{5}$ \quad (c) $\begin{cases}x=1\\y=2\end{cases}$ \quad (d) $x_1=7,x_2=-1$\\
2. $3<x\le5$\\
3.(a) $x_1=3,x_2=-\dfrac12$ \quad (b) $x_1=0,x_2=4$（验根：$x=0$应舍去，只有$x=4$）\\
4.(a) 54 \quad (b) 最多打8折 \quad (c) 甲$\dfrac{900}{11}$千米/时，乙$\dfrac{750}{11}$千米/时\\
5. 因为每次变形（去分母、平方、因式分解降次）最终都归结为一次或二次方程。

\sectiontitle{第七章}

\textbf{第18讲}

1.(a) $-1$ \quad (b) $-i$ \quad (c) 1 \quad (d) $-1$ \quad (e) $i$ \quad (f) 1\\
2.(a) 0 \quad (b) 1 \quad (c) 1\\
3.(a) 错 \quad (b) 错 \quad (c) 对 \quad (d) 对\\
4. 二次方程判别式<0时直接说无解，不影响公式的"完整性"。三次方程必须有实根，但公式却要处理负数开平方，所以逼出了新数。

\textbf{第19讲}

1.(a) $6-2i$ \quad (b) $-2-6i$ \quad (c) $8-i$ \quad (d) 5 \quad (e) $8+6i$ \quad (f) 2\\
2.(a) $3-4i$ \quad (b) $2+i$ \quad (c) $-5i$ \quad (d) $-1+2i$\\
3. $z+\bar{z}=2a$是实数；$z-\bar{z}=2bi$是纯虚数（或0）。\\
4. 略（复平面作图）\\
5.(a) $x=\pm i$ \quad (b) $x=\pm 2i$ \quad (c) $x=1\pm i$ \quad (d) $x=-1\pm 2i$\\
6. $\mathbb{N}$（自然数）→减法不够→$\mathbb{Z}$（整数）→除法不够→$\mathbb{Q}$（有理数）→开方不够→$\mathbb{R}$（实数）→负数开平方不够→$\mathbb{C}$（复数）
